%% ----------------------------------------------------------------------------
%%  §9.4  Integrazione delle funzioni razionali
%% ----------------------------------------------------------------------------
\section{Integrazione delle funzioni razionali}
\label{sec:razionali}

È sempre possibile, in linea di principio, calcolare per decomposizione in somma
l'integrale indefinito delle \emph{funzioni razionali}, cioè delle funzioni che sono
rapporto di due polinomi:
\begin{equation}\label{eq:funzione-razionale}
  \frac{f(x)}{g(x)}=\frac{a_{m}x^{m}+a_{m-1}x^{m-1}+\dots+a_{1}x+a_{0}}
  {b_{n}x^{n}+b_{n-1}x^{n-1}+\dots+b_{1}x+b_{0}}, \qquad a_{m},b_{n}\neq0 .
\end{equation}
Il procedimento si svolge in due tempi: prima ci si riconduce al caso in cui il grado del
numeratore è minore di quello del denominatore, poi si scompone la frazione in frazioni
«semplici», i cui integrali sono immediati.

\subsection{La divisione tra polinomi}

Se $m\ge n$, si esegue la \emph{divisione} tra i polinomi $f$ e $g$: come si dimostra in
algebra (per induzione sul grado di $f$), esistono e sono unici due polinomi $q$, il
\emph{quoziente}, e $r$, il \emph{resto}, con $r$ nullo oppure di grado minore di $n$, tali
che
\begin{equation}\label{eq:divisione-polinomi}
  f(x)=g(x)\,q(x)+r(x), \qquad\text{cioè}\qquad \frac{f(x)}{g(x)}=q(x)+\frac{r(x)}{g(x)} .
\end{equation}
Il polinomio $q$ si integra immediatamente, e resta da integrare la funzione razionale
$\frac rg$, nella quale il grado del numeratore è minore di quello del denominatore.

\begin{esempio}[Una divisione]\label{es:divisione-polinomi}
Dividiamo $f(x)=x^{5}-3x^{4}+x+3$ per $g(x)=x^{2}-1$, procedendo come nella divisione tra
numeri naturali: a ogni passo si divide il termine di grado massimo del dividendo
parziale per $x^{2}$, si moltiplica il risultato per $g$ e lo si sottrae:
\[
\begin{aligned}
  x^{5}-3x^{4}+x+3-x^{3}(x^{2}-1)&=-3x^{4}+x^{3}+x+3,\\
  -3x^{4}+x^{3}+x+3+3x^{2}(x^{2}-1)&=x^{3}-3x^{2}+x+3,\\
  x^{3}-3x^{2}+x+3-x(x^{2}-1)&=-3x^{2}+2x+3,\\
  -3x^{2}+2x+3+3(x^{2}-1)&=2x .
\end{aligned}
\]
Il quoziente è $q(x)=x^{3}-3x^{2}+x-3$ e il resto è $r(x)=2x$. Poiché $2x$ è la derivata
di $x^{2}-1$,
\[
\begin{aligned}
  \int\frac{x^{5}-3x^{4}+x+3}{x^{2}-1}\,dx&=\int(x^{3}-3x^{2}+x-3)\,dx
  +\int\frac{2x}{x^{2}-1}\,dx\\
  &=\frac{x^{4}}4-x^{3}+\frac{x^{2}}2-3x+\log\abs{x^{2}-1}+c .
\end{aligned}
\]
\end{esempio}

\subsection{Denominatore di secondo grado}

Consideriamo ora una funzione razionale $\frac{r(x)}{g(x)}$ con $g$ di secondo grado e $r$ di
grado al più uno. Raccogliendo il coefficiente direttore di $g$ ci si riduce al caso
\begin{equation}\label{eq:razionale-secondo-grado}
  \frac{mx+n}{x^{2}+px+q}, \qquad \Delta=p^{2}-4q ,
\end{equation}
e occorre distinguere tre casi, secondo che l'equazione $x^{2}+px+q=0$ abbia due radici
reali distinte, due radici coincidenti, o nessuna radice reale.

\begin{proposizione}[Integrali delle frazioni con denominatore di secondo grado]
\label{prop:razionali-secondo-grado}
Siano $m,n,p,q\in\R$ e $\Delta=p^{2}-4q$.
\begin{enumerate}[label=\textup{(\roman*)}]
  \item Se $\Delta>0$ e $x_{1}\neq x_{2}$ sono le radici di $x^{2}+px+q$, esistono e sono unici
  $A,B\in\R$ tali che
  \begin{equation}\label{eq:fratti-semplici-distinte}
  \begin{gathered}
    \frac{mx+n}{(x-x_{1})(x-x_{2})}=\frac A{x-x_{1}}+\frac B{x-x_{2}},\\
    A=\frac{mx_{1}+n}{x_{1}-x_{2}}, \qquad B=\frac{mx_{2}+n}{x_{2}-x_{1}},
  \end{gathered}
  \end{equation}
  e quindi l'integrale vale $A\log\abs{x-x_{1}}+B\log\abs{x-x_{2}}+c$.
  \item Se $\Delta=0$ e $x_{1}$ è la radice doppia,
  \begin{equation}\label{eq:fratti-semplici-doppia}
    \int\frac{mx+n}{(x-x_{1})^{2}}\,dx=m\log\abs{x-x_{1}}-\frac{mx_{1}+n}{x-x_{1}}+c .
  \end{equation}
  \item Se $\Delta<0$, posto $\delta=\frac{\sqrt{-\Delta}}2>0$,
  \begin{equation}\label{eq:fratti-semplici-complesse}
  \begin{split}
    \int\frac{mx+n}{x^{2}+px+q}\,dx&=\frac m2\log\bigl(x^{2}+px+q\bigr)\\
    &\quad+\frac{n-\frac{mp}2}{\delta}\,\operatorname{arctg}\frac{x+\frac p2}\delta+c .
  \end{split}
  \end{equation}
\end{enumerate}
\end{proposizione}

\begin{schema}
\Ipotesi{$\frac{mx+n}{x^{2}+px+q}$, con discriminante $\Delta$ positivo, nullo o negativo.}
\Tesi{le decomposizioni e le primitive \eqref{eq:fratti-semplici-distinte},
\eqref{eq:fratti-semplici-doppia}, \eqref{eq:fratti-semplici-complesse}, ciascuna valida in
ogni intervallo in cui il denominatore non si annulla.}
\end{schema}

\begin{strategia}
Nel primo caso si cercano $A$ e $B$ imponendo che la somma delle due frazioni, ridotta
allo stesso denominatore, abbia numeratore $mx+n$: è un'identità tra polinomi, che basta
valutare nelle radici per trovare $A$ e $B$. Nel secondo caso si scrive il numeratore come
$m(x-x_{1})+(mx_{1}+n)$. Nel terzo il denominatore non si scompone; si scrive il numeratore
come multiplo della derivata $2x+p$ del denominatore (che dà un logaritmo) più una
costante, e si completa il quadrato al denominatore, che diventa $(x+\frac p2)^{2}+\delta^{2}$
(che dà un arcotangente).
\end{strategia}

\begin{dimostrazione}
(i) Riducendo allo stesso denominatore, la \eqref{eq:fratti-semplici-distinte} equivale a
\[
  mx+n=A(x-x_{2})+B(x-x_{1}) \qquad\text{per ogni } x\neq x_{1},x_{2} .
\]
I due membri sono polinomi di grado al più uno che coincidono in infiniti punti, quindi
hanno gli stessi coefficienti e coincidono per ogni $x$ (la differenza è un polinomio
di grado al più uno con più di una radice, dunque nullo). Ponendo $x=x_{1}$ si ottiene
$mx_{1}+n=A(x_{1}-x_{2})$, e ponendo $x=x_{2}$ si ottiene $mx_{2}+n=B(x_{2}-x_{1})$: $A$ e $B$ sono
determinati univocamente come nella \eqref{eq:fratti-semplici-distinte}; con tali valori
l'identità è verificata, perché due polinomi di grado al più uno che coincidono in due
punti distinti sono uguali. L'integrale segue da $\int\frac{dx}{x-x_{i}}=\log\abs{x-x_{i}}+c$.

(ii) Si ha $mx+n=m(x-x_{1})+(mx_{1}+n)$, e quindi
$\frac{mx+n}{(x-x_{1})^{2}}=\frac m{x-x_{1}}+\frac{mx_{1}+n}{(x-x_{1})^{2}}$; poiché
$\int(x-x_{1})^{-2}dx=-(x-x_{1})^{-1}+c$, si ottiene la \eqref{eq:fratti-semplici-doppia}.

(iii) Scriviamo $mx+n=\frac m2(2x+p)+\bigl(n-\frac{mp}2\bigr)$. Il primo addendo dà, per la
seconda delle \eqref{eq:integrali-composti}, $\frac m2\log(x^{2}+px+q)$, senza valore
assoluto perché il denominatore, privo di radici reali, è sempre positivo. Per il
secondo, completiamo il quadrato:
\[
  x^{2}+px+q=\Bigl(x+\frac p2\Bigr)^{2}+q-\frac{p^{2}}4=\Bigl(x+\frac p2\Bigr)^{2}+\delta^{2}
  =\delta^{2}\Bigl[1+\Bigl(\frac{x+p/2}\delta\Bigr)^{2}\Bigr] ,
\]
con $\delta^{2}=q-\frac{p^{2}}4=-\frac\Delta4>0$. Per la \eqref{eq:argomento-lineare} con
$\alpha=\frac1\delta$,
\[
\begin{aligned}
  \int\frac{dx}{x^{2}+px+q}&=\frac1{\delta^{2}}\int\frac{dx}{1+\bigl(\frac{x+p/2}\delta\bigr)^{2}}
  =\frac1{\delta^{2}}\cdot\delta\,\operatorname{arctg}\frac{x+p/2}\delta+c\\
  &=\frac1\delta\operatorname{arctg}\frac{x+p/2}\delta+c ,
\end{aligned}
\]
e la \eqref{eq:fratti-semplici-complesse} segue per la linearità.
\end{dimostrazione}

\begin{esempio}[I tre casi]\label{es:tre-casi}
\begin{enumerate}[label=\textup{(\alph*)}]
  \item $\Delta>0$. Le radici di $x^{2}-x-2$ sono $x_{1}=-1$ e $x_{2}=2$. Per la
  \eqref{eq:fratti-semplici-distinte} con $m=1$, $n=7$: $A=\frac{-1+7}{-3}=-2$,
  $B=\frac{2+7}3=3$, e
  \[
  \begin{aligned}
    \int\frac{x+7}{x^{2}-x-2}\,dx&=\int\Bigl(\frac{-2}{x+1}+\frac3{x-2}\Bigr)dx\\
    &=-2\log\abs{x+1}+3\log\abs{x-2}+c .
  \end{aligned}
  \]
  \item $\Delta=0$. Per $\frac x{x^{2}+2x+1}=\frac x{(x+1)^{2}}$ si ha $x_{1}=-1$, $m=1$, $n=0$:
  \[
    \int\frac x{(x+1)^{2}}\,dx=\int\Bigl(\frac1{x+1}-\frac1{(x+1)^{2}}\Bigr)dx
    =\log\abs{x+1}+\frac1{x+1}+c .
  \]
  \item $\Delta<0$. Per $\frac{1-2x}{x^{2}+2x+5}$ si ha $p=2$, $q=5$, $\delta=2$, $m=-2$, $n=1$:
  \[
    \int\frac{1-2x}{x^{2}+2x+5}\,dx=-\log(x^{2}+2x+5)+\frac32\operatorname{arctg}\frac{x+1}2+c .
  \]
  Allo stesso modo, per $\frac x{x^{2}+x+1}$ ($p=q=1$, $\delta=\frac{\sqrt3}2$, $m=1$, $n=0$),
  \[
    \int\frac x{x^{2}+x+1}\,dx=\frac12\log(x^{2}+x+1)-\frac{\sqrt3}3\operatorname{arctg}
    \frac{2x+1}{\sqrt3}+c .
  \]
\end{enumerate}
\end{esempio}

\subsection{Denominatori di grado superiore}

Il metodo si estende a denominatori di grado qualsiasi, purché se ne conoscano
esplicitamente le radici. Si dimostra in algebra (è una conseguenza del teorema
fondamentale dell'algebra, ricordato nel \cref{cap:complementi}) che ogni polinomio
reale si scompone nel prodotto di fattori di primo grado $x-x_{i}$ e di fattori di secondo
grado $x^{2}+px+q$ con $\Delta<0$, eventualmente ripetuti; e che ogni funzione razionale con
il numeratore di grado minore del denominatore è somma di \emph{fratti semplici} del
tipo
\[
  \frac A{(x-x_{i})^{k}}, \qquad \frac{Bx+C}{(x^{2}+px+q)^{k}} ,
\]
uno per ogni potenza di ciascun fattore che compare nel denominatore. Le costanti si
determinano, come nella \cref{prop:razionali-secondo-grado}, imponendo l'identità dei
numeratori.

\begin{esempio}[Un denominatore di quarto grado]\label{es:quarto-grado}
Dividendo il numeratore per il denominatore,
$\frac{x^{5}-x+1}{x^{4}+x^{2}}=x-\frac{x^{3}+x-1}{x^{4}+x^{2}}$. Il denominatore è
$x^{2}(x^{2}+1)$, con il fattore $x$ ripetuto due volte e il fattore irriducibile $x^{2}+1$;
cerchiamo dunque
\[
  \frac{x^{3}+x-1}{x^{2}(x^{2}+1)}=\frac Ax+\frac B{x^{2}}+\frac{C+Dx}{x^{2}+1} .
\]
Riducendo allo stesso denominatore e uguagliando i coefficienti dei numeratori
($x^{3}+x-1=Ax(x^{2}+1)+B(x^{2}+1)+(C+Dx)x^{2}$) si trova $A=C=1$, $B=-1$, $D=0$, e quindi
\[
  \int\frac{x^{5}-x+1}{x^{4}+x^{2}}\,dx=\frac{x^{2}}2-\log\abs x-\frac1x-\operatorname{arctg}x+c .
\]
\end{esempio}

\begin{python}[Divisione, fratti semplici e integrazione con sympy]
import sympy as sp

x = sp.symbols("x", real=True)
q, r = sp.div(x**5 - 3*x**4 + x + 3, x**2 - 1, x)
print("quoziente:", q, "   resto:", r)
for f in [(x + 7) / (x**2 - x - 2), x / (x**2 + 2*x + 1),
          (1 - 2*x) / (x**2 + 2*x + 5),
          (x**5 - x + 1) / (x**4 + x**2)]:
    print("f =", f)
    print("    fratti semplici:", sp.apart(f, x))
    print("    integrale:      ", sp.integrate(f, x))
\end{python}
\begin{risultato}
quoziente: x**3 - 3*x**2 + x - 3    resto: 2*x
f = (x + 7)/(x**2 - x - 2)
    fratti semplici: -2/(x + 1) + 3/(x - 2)
    integrale:       3*log(x - 2) - 2*log(x + 1)
f = x/(x**2 + 2*x + 1)
    fratti semplici: 1/(x + 1) - 1/(x + 1)**2
    integrale:       log(x + 1) + 1/(x + 1)
f = (1 - 2*x)/(x**2 + 2*x + 5)
    fratti semplici: -(2*x - 1)/(x**2 + 2*x + 5)
    integrale:       -log(x**2 + 2*x + 5) + 3*atan(x/2 + 1/2)/2
f = (x**5 - x + 1)/(x**4 + x**2)
    fratti semplici: x - 1/(x**2 + 1) - 1/x + x**(-2)
    integrale:       x**2/2 - log(x) - atan(x) - 1/x
\end{risultato}

Il calcolo simbolico conferma le decomposizioni e le primitive (ancora una volta senza
valori assoluti nei logaritmi). Si noti però che il metodo richiede le radici del
denominatore: per un polinomio di grado elevato esse si calcolano, in generale, soltanto
numericamente, e la decomposizione in fratti semplici diventa un procedimento
approssimato.

\begin{notastorica}[Leibniz, Johann Bernoulli e i fratti semplici]
Il metodo dei fratti semplici fu pubblicato quasi contemporaneamente, nel 1702, da
Leibniz e da Johann Bernoulli sugli \emph{Acta Eruditorum}. Leibniz incontrò una
difficoltà che lo portò ad affermare, erroneamente, che non tutte le funzioni razionali
si possono integrare con logaritmi e arcotangenti: non riusciva a scomporre $x^{4}+1$ in
fattori reali di secondo grado. A posteriori l'errore è evidente:
$x^{4}+1=(x^{2}+\sqrt2\,x+1)(x^{2}-\sqrt2\,x+1)$, come si verifica moltiplicando. Il
problema fu risolto in generale soltanto dal teorema fondamentale dell'algebra, che nacque
dunque come questione di calcolo integrale: il primo tentativo di dimostrarlo, di Jean Le
Rond d'Alembert (1746), compare in una memoria intitolata \emph{Recherches sur le calcul
intégral}. La prima dimostrazione sostanzialmente completa, pubblicata da Gauss nel 1799,
conteneva una lacuna topologica colmata da Alexander Ostrowski nel 1920. In Inghilterra la
scuola di Newton seguiva un'altra strada: sviluppava in serie di potenze la funzione
integranda e integrava termine a termine, senza bisogno di scomporre il denominatore. Il
metodo era efficace per i calcoli numerici, ma non dava le formule chiuse in termini di
logaritmi e arcotangenti cercate dai continentali.
\end{notastorica}
